Language-model agents are moving from answering questions to doing research. Systems such as ChemCrow~\citep{bran2024chemcrow} and
Biomni~\citep{huang2026biomni} plan multistep analyses, query scientific databases and compute on what they retrieve. Their answers
can only be as good as the records they read, and those records contain errors. Misannotation is widespread in several major
protein databases~\citep{schnoes2009annotation}, and even curated resources revise their entries: UniProt archives every earlier
version of every entry~\citep{leinonen2006unisave}, and of the 300 reviewed human proteins in our benchmark, 167 have at least one
superseded sequence. An agent that relays a database value inherits its error unless it notices the error and acts on it.

How language models handle wrong information in their context has been studied mainly as a conflict between the context and the
model's own knowledge~\citep{longpre2021entity,xu2024conflicts}. Models often adopt incorrect retrieved content, less often the
further it departs from the truth~\citep{wu2024clasheval}, and they accept counter-evidence that is coherent and convincing even when
it contradicts their memory~\citep{xie2024adaptive}. Concurrent work shows that tool-using agents adopt corrupted tool returns and
may question a return, or even state the correct value, while passing the corrupted content on~\citep{yang2026overreliance}. These
studies corrupt passages of text or single returned values. Scientific database records are different. A value sits among related
fields that can confirm or contradict it, and the record often contains what is needed to derive the value independently: a
protein's sequence fixes its length and mass, and a compound's formula fixes its weight. Whether agents use this structure, whether
errors that occurred in real databases behave like constructed ones, and which corrections actually change what an agent reports,
are open questions.

We studied them in five steps. We ran a deployed biomedical agent against live database APIs and altered the responses in transit.
We built a controlled benchmark of protein and small-molecule questions and ran ten models from five developers. We manipulated the
internal consistency of the records and the derivation routes they offer. We presented errors that UniProt itself later corrected.
And we tested corrections: checks of each record against its own derivations, an independent second source, and three ways of
presenting a correction to the agent. Where we made predictions, we fixed them in dated repository commits before running the experiments, deposited the deployment hypotheses on OSF, and report the ones that failed.

Across these experiments, agents noticed wrong values far more often than they acted on them. Whether a wrong value was relayed
depended less on how plausible it was than on whether the rest of the record agreed with it: adding one agreeing field raised relay,
and making the record disagree with itself lowered it, with the wrong value unchanged. A derivation route let agents detect an error,
and often compute the correct value, without changing what they answered. Errors that UniProt later corrected, each internally
consistent, were relayed in 90--100\% of runs, whatever the size of the revision. An independent second source removed most of
these errors; correcting the record silently removed all of them in our tests, whereas showing the agent the original value invited
it back.
